\section{Polynomial matching consequences before Q enumeration}
\label{sec:matching-support}

The complete-Q support route of Section~\ref{sec:feasible-span} still pays
for ray enumeration in the original matching dimension. This continuation
adds a polynomial sufficient preprocessor before that enumeration. It
certifies coordinate zeros from exact matching equations, then recompiles
the source using only the retained types. Its checker verifies source-row
combinations directly, without trusting a producer's forest, reduced matrix
or inferred support. The complete Q and standard fallbacks remain available
when the inexpensive implications do not find every forced zero.

\subsection{One-sided equations imply zeros}

Suppose earlier verified steps have forced coordinates in $F$ to vanish.
For an actual matching consequence $w\cdot q=0$, assume
$w_j\ge0$ for every allowed coordinate outside $F$. Then every still-active
coordinate with $w_j>0$ is zero in every feasible vector: the equality is
a sum of nonnegative terms, and the already forced terms vanish.
A row whose active coefficients are all negative is multiplied by $-1$.
Rows with both signs are left for the complete search.

\begin{proposition}[Sound polynomial sign propagation]
Starting with $F=\varnothing$, repeatedly applying the preceding implication
to a fixed rational row family preserves every feasible standard vector.
At most $k$ successful steps occur when $k$ types are allowed. Removing the
proved-zero coordinates and restoring them as zeros gives a linear
isomorphism between the original and restricted standard cones.
\end{proposition}
\begin{proof}
Induct on the steps. Every preceding zero is valid by the induction
hypothesis, and the remaining summands in the new matching consequence
are nonnegative. A positive coefficient therefore forces its coordinate
to zero. Each successful step removes at least one new coordinate, giving
the step bound. Every original feasible standard vector already has the
removed coordinates zero, while reinstating them in a restricted feasible
vector preserves all original matching and nonnegativity equations. The
resulting isomorphism preserves the non-link standard rays and actual
source coordinates.
\end{proof}

This is the preceding single-row implication procedure.
Section~\ref{sec:matching-pairs} strengthens it using two-row spans;
neither rule is a complete maximal-support
LP solver. For example the two RREF rows
\[
 (1,0,-1,2),\qquad(0,1,2,-1)
\]
have mixed signs and yield no step. Their sum is $(1,1,1,1)$, so their
nonnegative matching cone is actually zero. This counterexample separates
soundness from completeness for general row systems. In particular a
mixed-sign fixed point does not certify positive feasibility, maximal
support or absence of an essential disc.

\subsection{Recover useful rows cheaply, then compile source provenance}

The existing canonical nullspace basis $K$ has one unit free coordinate
per row. For each pivot coordinate $p$, its matching RREF equation is
\[
             q_p-\sum_a K_{a,p}q_{f_a}=0.
\]
Recovering these rows from the already computed basis costs $O(k^2)$ scalar
work. Repeated sign propagation on at most $k$ such rows costs $O(k^3)$
scalar inspections; rational sign comparisons and coefficient bits are
charged. If no coordinate is forced, no source provenance forest is built.

When useful zeros are found, the compiler reconstructs triangle incidence
from the actual source matching equations. A directed edge with equation
$t_a-t_b+L(q)=0$ has label $L$ and a unit source-equation coefficient.
A rooted forest stores both its Q potential $P_c$ and the source-row
combination $E_c$ satisfying
\[
                   E_c=t_{\mathrm{root}}-t_c+P_c(q).
\]
For an edge $a\to b$, the combination $E_a+E_{a\to b}-E_b$ cancels every
triangle coefficient and gives $P_a+L-P_b$. Source equations with no
triangle terms supply direct Q constraints. Exact elimination pivots only
on the allowed Q columns; appended source-row coefficient columns record
all the same row operations. The resulting Q RREF rows must equal the rows
recovered from the native kernel, and their sign closure must agree.
These exact equalities detect a provenance/gauge disagreement before any
restricted geometry is used.

The forest, path combinations and augmented elimination have polynomial
source-size and bit cost. They are not charged as constant-time preparation:
a direct implementation can retain quadratic source path data. Coefficient
bit growth is bounded by the usual minor bounds for exact rational
elimination. This additional work is compiled only when the cheap recovered
rows already show a valid reduction.

\subsection{Independent linear replay licenses source recompilation}

The optional \code{matching\_support\_certificate} records the original
source digest and allowed type list. Its schema is
\code{normal-sector-matching-support-v1}. Each step contains sorted source-row indices, exact
numerator/denominator coefficients and the newly forced coordinates.
The independent checker prepares the source anew and forms each stated
combination of its actual matching equations. It checks cancellation of
\emph{all} triangle coefficients, nonnegative coefficients on every
currently unforced allowed Q coordinate, and exact agreement between the
claimed new coordinates and those with strictly positive coefficients.
Previously forced coefficients need no sign condition because their
coordinates have already been authenticated as zero. Nonselected Q types
are zero by the original allowed-sector definition.

Malformed source indices, Boolean indices, zero or invalid denominators,
wrong source/support identities, repeated indices, unsupported nesting,
nonzero triangle remnants and incomplete coordinate claims are rejected.
The checker imports neither the producer's forest nor its RREF/closure
helpers. Valid coefficients can be represented by the existing large-integer
codec; replay uses exact rational arithmetic and no numerical tolerance.

Automatic standard discovery above nullity three runs this preprocessor
before a generic Q phase. A proper reduction rebuilds the source kernel and
resumes automatic discovery in its smaller matching space, with the same
candidate allowance. Low-dimensional progressive corner/envelope/planar
branches can therefore run without the original-dimensional Q scan. An
ordinary positive witness retains the original allowed list and full source
coordinates; its disc proof suffices independently of this preprocessing.

Negative parent certificates also retain the original support and coordinate
width. Their checker authenticates the matching implications first and
reconstructs the retained model independently. If later complete Q support
removes more types, the existing Q support proof concerns that intermediate
retained sector; replay checks that precise scope, derives the next support
and restores all coordinates to the original parent width. The matching
and Q support proofs can therefore compose without narrowing the final
statement. Removing all support annotations still leaves an original full-sector
proof that the legacy checker can reconstruct directly.

\subsection{The raw-dimensional Q cost can now disappear}

Let $r$ be the matching nullity after the sufficient polynomial preprocessing,
and $k'$ its retained support size. The complete ray-geometry search has
\[
 \operatorname{poly}(T+k+B)+
                    \operatorname{poly}(T+B)(1+k')^{O(r)}
\]
bit cost, including the Q prelude, standard fallback, source recompilation
and source-equation certificate reading. The old raw nullity need not
appear in the enumerative term. Thus logarithmic $r$ gives a conditional
quasi-polynomial geometry bound even if the original matching dimension
was larger. At $r\le3$ the native low-dimensional geometry applies directly.
Component queries, essentiality and their proof costs remain separate.

The retained matching dimension can exceed the feasible cone's actual span
when the sign procedure misses zeros; no maximality is asserted from its
fixed point alone. Neither logarithmic $r$ on arbitrary knot exteriors nor
a bounded complete family of sectors or centres has been proved. This
implements a polynomial sufficient operation toward that goal, without a
new general recognition bound or an implementation of the full LP theorem.

\subsection{Source controls and complete-query measurements}
All 1,433 maintained tests pass in 259.100 seconds. The 34 focused tests
pass in 13.277 seconds. They cover source-row planar admission before any
generic Q scan, exact rational signs, fourteen source-provenance mutations,
legacy full-source replay and composition of a partial matching proof with
the complete Q-support fallback. The mixed-row counterexample is tested as
an explicit non-completeness boundary. The preceding complete-Q regression
route is still exercised with matching preprocessing disabled internally.

The 47.730-second audit repeats all 188 retained higher-nullity source
sectors against release \code{bc98f99a9}. Every complete status agrees.
For this entire selected corpus, the polynomial recovered-row closure finds
all 126 oracle forced-zero coordinates in 65 sectors; this observation is
not promoted to a general maximal-support theorem. Fresh Regina standard
enumerations for all three involved source triangulations reproduce every
compatible full-coordinate ray list. The full support oracle again gives
22 transitions from dimension four to three, 41 reductions that stay at
four and two transitions from five to four.

All 65 affected automatic searches now recompile before their generic Q
phase. Four return an ordinary verified essential-disc witness and 61
return negative proofs with the source-equation annotation. The 22 sectors
that drop to dimension three enter the progressive planar producer without
the original-dimensional Q prelude. A returned ordinary positive witness
requires no additional support annotation. Every negative proof retains
the original allowed types and full coordinate width. Independent review
replays all 61 annotated proofs and all 31 saved diagram witnesses with the
new forcing, forest, Q and ray producers disabled. Frozen original
full-source replay also accepts representative unannotated proofs at
effective dimensions three and four.

The complete-query timing driver freezes both the previous discovery code
and its independent checker. Each arm charges source construction, its
complete configured search, every component query and proof replay, and
full serialized output. Five randomized paired rounds with identical-baseline
A/A controls complete 100 measured calls and twenty warmups. Every status
agrees across arms, and each arm's certificate digest is stable across
its repetitions.

\begin{center}
\small
\begin{tabular}{lrrrrr}
\toprule
Source & previous ms & current ms & old/new & old A/A & new A/A\\
\midrule
\code{cap\_3\_5\_2-4-to-3} & 40.904 & 25.583 & 1.614 & 1.002 & 0.982\\
\code{cap\_3\_5\_3-4-to-4} & 80.217 & 63.367 & 1.255 & 0.999 & 1.002\\
\code{cap\_3\_5\_3-5-to-4} & 142.674 & 122.678 & 1.166 & 1.002 & 0.983\\
\code{cap\_1\_2\_3-4-to-4} & 41.597 & 41.056 & 1.023 & 0.990 & 1.000\\
\code{fibonacci\_lst\_09-4-to-4} & 63.082 & 61.222 & 1.040 & 1.027 & 1.013\\
\bottomrule
\end{tabular}
\end{center}


The planar-admission query has paired old/new ratio 1.614 over the already
improved preceding release. Its begun candidates fall from nine to five;
the original four Q-basis candidates are no longer required. Serialized
result size falls from 6,569 to 4,451 bytes. Its A/A controls are about 1.002
and 0.982. The affected same-dimension-four query gains 1.255 with the same
124 candidates; its linear source proof avoids the preceding complete-Q
support reconstruction and reduces serialized size from 10,468 to 7,822
bytes. A/A controls are about 0.999 and 1.002.

The dimension-five-to-four query gains 1.166, with 369 versus 368 candidates
and serialized size 11,820 versus 8,889 bytes. Its A/A controls are about
1.002 and 0.983. The one-candidate change is not the whole gain: direct
linear implication replay replaces authentication of an original-dimensional
Q list and its component proofs. Source recompilation and provenance
preparation are included rather than treated as cached or zero-cost.

The full-support positive and nonpositive Q controls have ratios 1.023 and
1.040. Their A/A observations include ratios about 0.990 and 1.027; no
reliable improvement is inferred on those unchanged paths. Recovering and
scanning reduced rows also has a cost when no type is removed. These are
host measurements of supplied-source queries, not a general speed guarantee
or a claimed complete knot-recognition gain.

Runtime release \code{4c75b66de}, tests, the complete source audit,
fresh Regina controls and paired timings are retained under
\code{data/matching-support-*}.

The driver is the module
\code{normal\_orbit\_research.}\allowbreak\code{matching\_support},
run from \code{fast/} with \code{audit} or \code{benchmark}.
The audit accepts \code{--fresh-regina}.
Publication review checks 612 runtime pins, 609 preceding-release pins,
source oracles, all timing outcomes, and the producer-disabled source and
saved-diagram replays. The meaningful asymptotic change is sufficient
polynomial preprocessing before enumeration; a complete polynomial support
solver, universal logarithmic remaining dimension and general recognition
coverage remain unimplemented or unproved.

